\section{Consequences for exact algorithms and convexification}\label{s6:synthesis}

The $\ETR$ reductions realize arbitrary bounded
algebraic constraints within physical mixing equations. The one- and two-pool
$\NP$-hardness reductions retain only a few mixing variables but couple
them to a large linear network. The structural and contract algorithms
succeed when that network can be eliminated while retaining all balances
and the objective. These statements classify decision or optimization;
they do not predict the running time of a particular numerical solver.
This section states the representation consequences of three supporting
results, proved in Appendices~\ref{app:geometry}--\ref{app:convexification},
and the main question left open.

\subsection{Large response descriptions can coexist with exact optimization}
A pool-free path is already enough for an exponential parametric response.
\Cref{s6:price-response} constructs a rational physical family with
$O(n)$ nodes, one upper-bounded attribute, only two distinct input
qualities, zero flow lower bounds, and positive product revenues.
Changing only the final product's price exposes $2^n$ distinct unique
optimal flows and $2^n$ affine value regimes. The construction has
polynomial encoding length; its exact penalties have large, polynomially
encoded revenues. Every fixed-price problem is an LP.
\Cref{s6:response-size} proves a matching lower bound for explicit
piecewise-affine tables and for the total polynomial degree of formulas
in price and value alone. It leaves compact extended descriptions and
implicit evaluation available.

A single two-feed, two-outlet pool attached to the contracted path yields
$2^n$ isolated global optima with ordinary nonnegative input costs and
product revenues at most two (\cref{s6:vertex-forcing}). A rational price
perturbation yields $2^n$ strict local maxima with distinct values
(\cref{s6:local}); no tolerance or local-solver complexity bound follows.
Nevertheless, the full feasible set of this family has a polynomial-size
LP convex hull whose vertices are physically feasible
(\cref{s6:lp-hull}), so an optimal LP vertex is an exact physical
optimizer for every linear objective. The same
family requires exactly $n$ integer coordinates in any convex lift of
its finite optimal set (\cref{s6:integer-lift}). The represented set is
material: a convex hull preserves linear optimization but does not
represent the original nonconvex feasible or optimal set exactly.

These constructions realize an established exponential Klee--Minty
shadow \cite{s6:gartner2013} by physical blending constraints. Their role
is to rule out enumeration arguments, not to establish hardness on
paths. Indeed, the rank-one concave objective used in the construction
can be minimized over the bare path by a linear program. The abstract
hardness in \cref{s6:slab-hardness} requires an additional dense weighted
slab. No degree-two physical realization of that slab is proved here.

\subsection{Cost restrictions and uniform hull restrictions are different}
An isolated pool block records a nonnegative rank-one matrix with
interval row and column margins. General matrix costs make its rational
threshold problem strongly $\NP$-complete even when every lower margin
bound is zero and every upper bound is one
(\cref{app:margin-hardness,app:unit-hardness}). Every rational yes threshold
still has a polynomial-size rational matrix witness, and some exact
optimizer always lies in one quadratic number field
(\cref{app:margin-certificates}). Thus threshold witnesses can be simpler
than exact optimizers. These statements concern arbitrary
source-to-terminal costs on an isolated block; with ordinary additive
costs the quality-free margin problem is an LP, so the reduction is not
a physical pooling hardness proof.

The useful cost parameter is its rank after subtracting arbitrary
additive row and column costs. For each fixed value of this
\emph{interaction rank}, \cref{s6:rank-theorem} gives polynomial exact
optimization and an optimizer in one quadratic field. This yields a
fixed-attribute Lagrangian oracle after terminal quality inequalities
are dualized (\cref{s6:oracle}). The theorem is an oracle for the stated
isolated block; it implies neither zero duality gap nor an algorithm
for coupled physical pools.

The corresponding convex hull can still require large uniform
representations. Its unit-margin version contains an explicit face
isomorphic to the correlation polytope (\cref{app:correlation-face}).
Established correlation-polytope bounds then imply exponentially many
PSD blocks of any fixed order, exponential total second-order-cone
dimension, and superpolynomial unrestricted PSD order for an exact lift
(\cref{app:exact-conic}). A quantitative face estimate extends these
conclusions to uniform outer approximations: at the stated
inverse-quadratic total entrywise $\ell_1$ error, LP lifts need exponentially many
inequalities and SDP lifts need superpolynomial total matrix order
(\cref{app:approx-lp,app:approx-sdp}). These bounds allow arbitrary real
coefficients and do not assume that the lift is efficiently constructed.
They concern all bounded linear costs simultaneously, so they do not
exclude a useful relaxation for one objective or a fixed-rank cost family.

\citet[Section~3.1, after Theorem~4]
{dey2020-convexifications-of-rank-one-based} conjecture hardness of linear
optimization with simultaneous row and column bounds; the reduction in
Appendix~\ref{app:costs} proves that conjecture for this exact set.
\citet[Theorem~2]{s6:jalilian2026} prove second-order-cone
representability with row, column and total bounds by a generally
exponential construction; our size lower bounds are compatible with it
and say nothing about the empirical performance of their compact
relaxations. The correlation lower-bound
machinery is due to \citet{s6:fawzi2013,s6:lrs-full,s6:bfps2015}; the
proofs here supply the rank-one face and the quantitative transfer.

\subsection{What the parameter boundaries leave open}\label{s6:open}
The main remaining structural question is standard pooling with a fixed
number of pools, fixed affine input-quality rank, degree-two bypass
graphs, and unboundedly many pool attachments. General flow and quality
intervals are allowed; feasibility and dense-profit optimization should
be treated separately. Two issues remain: local endpoint relations
depend on unknown pool-quality parameters, and elimination must retain
pool mass and quality aggregates as well as any objective. For one
scalar parameter and general planar path relations the bound is
quasipolynomial (\cref{s5:quasipoly-path}); several parameters with
such relations remain open. The endpoint/objective scope is explained after
\cref{s4:degree-boundary},
and \cref{s6:slab-hardness} shows why a fixed aggregate count alone
does not suffice in an abstract system.

A shared fixed-quality helper pool does not automatically convert the
degree-three copy network into a degree-two bypass network. Each original
midpoint source has its own private supply equation; a shared helper
preserves only their total. For example, in two full midpoint-copy
pairs with endpoint qualities zero and two, midpoint quality one,
and demand four at each copy product, choose both endpoint-pair signals
zero in the first pair and two in the second. The required four midpoint
flows are $4,4,0,0$, totaling eight. A common helper with throughput eight
can supply them, but the original private midpoint supplies would demand
four in each pair. The local replacement therefore loses a required
relation, regardless of whether additional constraints elsewhere exclude
this particular assignment. The coordinate-port obstruction in
\cref{s6:ports} gives a separate reason arbitrary degree-three LP circuits
cannot simply be replaced by degree-two pool-free gadgets.

Several restrictions supply the missing aggregate information: a fixed
number of attachments (\cref{s4:bounded-attachments}), exact supplies
with a fixed number of receiving products (\cref{s5:fixed-products}),
full contracts or a fixed number of exceptions
(\cref{s5:quadratic-field,s5:arbitrary-exceptions,s5:restrictive-capacity}),
and two full input vectors (\cref{s5:two-vector-theorem}).
Table~\ref{s5:scope-table} gives their scopes. They do not combine
automatically: arbitrary exceptions with a restrictive common capacity,
or dense arc costs, require an additional argument.

The unrestricted one-pool, one-quality bypass threshold question is
settled by \cref{s3:base-bypass,s3:upper-only,s3:constant-thm}. Its
tractable predecessors retain material hypotheses: subdividing bypasses
into artificial pools changes pool count, and dropping source capacities
changes shared resource coupling. The remaining degree-two question
imposes restrictions that the hardness constructions do not meet.
Similarly, the two-quality response construction is an LP with interval
product specifications after penalties; it does not contradict the
exact-product-contract feasibility theorem.
